





\section{Lessons Learned from employing LLMs for query rewriting}
\label{gen:sec:insight}

We conducted a comprehensive experiment to study the applicability of LLMs for  rewriting of a variety of queries  and improving their performance. Specifically, we carefully reviewed each rewrite suggested by LLMs and evaluated it for correctness and latency. We paid particular attention to discerning which problems could be effectively addressed and which could not, along with identifying conditions under which the latter could be transformed into the former.  Our study not only provided valuable insight into the capabilities of LLMs and the challenges associated with using it for query rewriting but also contributed to the development and refinement of \genrewrite that is presented in \S\ref{gen:sec:approach}.
Below we summarize the main lessons that have informed \genrewrite's design decisions. 


\begin{comment}
    

\genhead{Lesson 1: LLMs are quite effective on certain workloads even with a simple prompt} Reasoning the performance issues of queries and ensuring correct rewriting is a non-trivial task. Surprisingly, even with the simple prompt shown in Figure~\ref{gen:fig:basic_prompt}, LLMs, by understanding the semantics inherent in the query, can find many rewrite opportunities that significantly improve performance but are ignored by rewrite rules or even human experts. \baseline works best on queries with low complexity and high frequency in LLMs' training corpus. A prime example is the Leetcode workload~\cite{gen_leetcode}. For instance, consider LeetCode problem \#176, wherein the goal is to find the second highest salary from the Employee table. \query1 in Figure~\ref{gen:fig:sql-comparison} is the human-written solution for this problem (submitted by LeetCode participants) while \query2 emerges as a rewrite generated by LLMs. \query2 is approximately 2000x faster than \query1 when executed on a database randomly populated with one million rows. 
The optimization process involves a deep understanding of the query intent through reasoning: (1) identifying that the goal of \query1 is to find the second-highest salary, (2) realizing that there is an alternative way to achieve the same thing, i.e., first find the maximum salary using a subquery and then compare all salaries with this maximum salary to find the second maximum, and finally (3) recognizing that \query2 achieves performance improvement by reducing the number of comparisons required. 

\end{comment}
% ~\cite{CiteSomePaperThatHasSaidSomethingSimilar}

% While LLMs have showcased remarkable capabilities in understanding the semantics of data, performing reasoning, and generating code, it is well-known that their capacity is not limitless. 
% Similarly, as queries become more intricate, the LLM's ability to identify rewriting opportunities diminishes. 





\genhead{Lesson 1: LLMs are effective on simple workloads but struggle with complex rewrites} While LLMs can identify effective rewrites for simpler queries using only basic prompts, their ability to recognize and apply meaningful transformations diminishes as query complexity increases.
Although it remains possible for LLMs to discover rewrite opportunities, the decreasing likelihood necessitates more iterations of LLM runs, resulting in significantly higher costs. In \S\ref{gen:sec:motivation}, we demonstrated how Q11 in TPC-DS workload could be optimized into a more efficient form by eliminating the unnecessary \sqlunionallcolor. Similar inefficiencies exist in other TPC-DS queries, such as Q4 and Q74. Among them, Q4 stands out as the most complex: it has the highest token count (1,170 vs. 723 in Q11 and 591 in Q74), the greatest number of \sqlunionallcolor segments (3 vs. 2 for the others), and the most joins (5 vs. 3 for the others). Despite ten attempts using the prompt shown in Figure~\ref{gen:fig:basic_prompt} with GPT-4, none of the rewrites for Q4 suggested removing the \sqlunionallcolor, a transformation that would have significantly reduced execution time. In contrast, 4 out of 10 attempts for Q11 and 6 out of 10 for Q74 moved in the correct rewriting direction. Note that a correct rewriting direction, such as the suggestion to eliminate unnecessary \sqlunionallcolor in Q74's rewrites, does not guarantee an error-free or equivalent rewrite to the original query. The inherent complexity of such queries makes it much more difficult for LLMs to generate fully correct rewrites, even when the high-level transformation logic is sound.


\begin{comment}
    


Q4 and Q74 from the TPC-DS benchmark present similar issues. Q4, in particular, is the more complex one, with the highest token count (1,170 for Q4, compared to 723 for Q11 and 591 for Q74), the most segments combined with UNION ALL (3, versus 2 in the others), and the highest number of joins (5 for Q4 against 3 for the others). Despite multiple attempts (10 runs with the prompt presented in Figure~\ref{gen:fig:basic_prompt} with GPT-4), none of the rewrites for Q4 recommended removing \sqlunionallcolor, a modification that could markedly reduce execution time. Conversely, 4 of 10 attempts for Q11 and 6 of 10 for Q74 moved towards a correct rewriting direction. This observation raises the question of whether we can bridge the gap between queries that LLMs can optimize and those LLMs cannot optimize through knowledge transfer, thereby enabling LLMs to optimize a broader range of queries, particularly those that are highly complex. This is the main motivation behind our notion of NLR2s.

Note that a correct rewriting direction, such as the suggestion to eliminate unnecessary \sqlunionallcolor in Q74's rewrites, does not guarantee an error-free or equivalent rewrite to the original query. The complexity of such queries exacerbates the challenge of manipulating their elements without introducing errors.


\genhead{Lesson 2: LLMs can optimize queries that they previously struggled to optimize when provided with suitable rewrite hints in the prompt}
As zero-shot prompting is not universally effective for query rewriting task,  it becomes essential to incorporate precise guidance within the prompts. Drawing inspiration from the few-shot prompting used in Text-to-SQL tasks~\cite{sun2306sql}, one approach is to provide LLMs with examples that include both the input query and its optimized version. Yet, an alternative strategy also holds promise. Consider the scenario where a human expert analyzes two queries that can be optimized using the same idea, for example, Q4 and Q74. The insight coming to one's mind is that both queries can benefit from ``avoid the unnecessary \sqlunionallcolor''. In fact, integrating this insight as a hint in the prompt consistently leads LLMs to suggest removing \sqlunionallcolor across different attempts. This latter method offers two benefits over the few-shot prompting approach: (1) it requires fewer tokens, thereby reducing costs and accelerating the rewrite generation process, and (2) the hints are easily interpretable by humans, facilitating quicker debugging. However, due to the time constraints of human experts, manually analyzing each query and suggesting beneficial rewrite hints is impractical. Therefore, there is a need for an automated method to identify rewrite hints tailored for each query.

\end{comment}

\genhead{Lesson 2: LLMs can optimize queries that they previously struggled to optimize when provided with suitable rewrite hints in the prompt}
As zero-shot prompting is not universally effective for query rewriting task,  it becomes essential to incorporate precise guidance within the prompts. Drawing inspiration from the few-shot prompting used in Text-to-SQL tasks~\cite{sun2306sql}, one approach is to provide LLMs with examples that include both the input query and its optimized version. Yet, an alternative strategy ---using concise, human-readable rewrite hints--- also holds promise.
For example, one may intuitively recognize that both queries can benefit from ``avoid the unnecessary \sqlunionallcolor''. In fact, integrating this insight as a hint in the prompt consistently leads LLMs to suggest removing \sqlunionallcolor across different attempts. This latter method offers two benefits over the few-shot prompting approach: (1) it requires fewer tokens, thereby reducing costs and accelerating the rewrite generation process, and (2) the hints are easily interpretable by humans, facilitating quicker debugging. However, due to the time constraints of human experts, manually analyzing each query and suggesting beneficial rewrite hints is impractical. Therefore, there is a need for an automated method to identify rewrite hints tailored for each query.



\begin{figure}[t]
\inv
  \centering
  \fbox{
    \begin{minipage}{0.8\linewidth}
        % This is a block of text inside a box in a figure in \LaTeX.
        % \textcolor{blue}{[input query]}\\
        \colorbox{gray!10}{[input query]}\\
      Rewrite this query to improve performance.
    \end{minipage}
  }
  \inv
  \captionsetup{font=small,name=Figure}
  \caption{The baseline approach for LLM-based query rewriting (\baseline)}
  \label{gen:fig:basic_prompt}
 \inv
\end{figure}

% but overlooked by LLMs in the zero-shot setting.





% We have underlined the importance of providing proper guidance and designing suitable prompt earlier in this section.   We observed that the zero-shot method (for example, the prompt presented in Figure~\ref{gen:fig:basic_prompt}) is not universally effective. Directly applying the few-shot learning approach from the Text-to-SQL problem is not effective, as one rewriting strategy does not fit all, and providing query pairs with rewrites irrelevant to the input query may confuse LLMs, negatively affecting the rewriting quality. Instead of supplying example query pairs, database experts can manually analyze the input query, summarize beneficial rewrites or transformations in natural language, and incorporate them into LLMs as hints. For instance, by including ``Replace \sqlunionallcolor with separate Common Table Expressions (CTEs)'' as a hint in the prompt for rewriting TPC-DS Q4, the success rate of recommending the removal of \sqlunionallcolor increases from 0\% to 100\%. However, due to the time constraints of human experts, manually analyzing each query and suggesting beneficial rewrite directions is impractical. Therefore, there is a need for an automated method to identify rewrite hints tailored for each query but overlooked by LLMs in the zero-shot setting.

% summarize beneficial rewrites or transformations in natural language


\begin{comment}
    Modern LLMs are equipped with advanced reasoning capability and 
broad general knowledge, which is essential in generating equivalent and more efficient rewrites. Advanced reasoning allows an LLM to look beyond superficial query structure, identify hidden redundancies, and streamline logical steps that a purely rule-based optimizer might overlook. \sqlunionallcolor and filtering operations canceling each other is a good example. LLMs are trained on large, diverse datasets, giving them exposure to a wide array of patterns, best practices, and domain-specific contexts. This breadth of information can surface rewriting strategies that might not exist in a limited, rule-based system. For example, \query1 in Figure~\ref{gen:fig:sql-comparison} is the human-written solution for this problem (submitted by LeetCode participants) while \query2 emerges as a rewrite generated by LLMs. \query2 is approximately 2000x faster than \query1 when executed on a database randomly populated with one million rows. The optimization process involves a deep understanding of the query intent through reasoning: (1) identifying that the goal of \query1 is to find the second-highest salary, (2) realizing that there is an alternative way to achieve the same thing, i.e., first find the maximum salary using a subquery and then compare all salaries with this maximum salary to find the second maximum, and finally (3) recognizing that \query2 achieves performance improvement by reducing the number of comparisons required. 
\end{comment}